\section{Threats to Validity}
\label{sec:ccd-threats}
\subsection{Internal Validity}
Inputs, quantization, tie rules, center order and owner/key pins must agree
across methods. Fresh public replay and independent integer costs check
these conditions in the retained runs. Resource stops and failed attempts
are preserved. In particular, the interrupted PCA run at $n=128$ is not a completed
baseline measurement. Proof-caching tests distinguish reusable old evidence
from newly generated request proofs.

The proof timing pilots have one execution per fit. Process loading, host
activity and the Windows/WSL environment can affect elapsed time. The final
text rerun is not paired with a newly timed Full arm. Earlier component
allocations are reconstructed from measured stages and are not new hot-RPC
experiments. The cost-aware policy was revised after inspecting the first
text result; it is engineering validation rather than a preregistered
independent confirmation. The later equal-routing gate has a newly measured
Full arm, and its seed-557 results must not be confused with that earlier
unpaired seed-521 rerun. Its three-fit numerical results do not create
three independent proof measurements. The RaBitQ and warm follow-ups were
registered after the first native outcomes; balanced accuracy is a post-hoc
descriptive metric. All sources have been inspected during development.

Implementation tests are evidence about specific failure cases. Accepted
decisions still depend on binding commitments, sound arguments, valid
integer ranges, complete coverage and an independently issued post-commit
challenge. The benchmark simulates those roles locally. Production setup
governance, authenticated challenge delivery and an external protocol audit
remain separate requirements.

\subsection{External Validity}
The current sources, fitted models, center counts and proof catalogs cover
a limited range of learning workloads. Image and text requests reuse their
training streams, references and audit catalogs; increasing query count does
not create additional independent learning trials. The small new pilots
cannot determine performance for arbitrary future centers, large prototype
banks, constantly changing references or different encoder distributions.

Native scalability and synthetic geometry do not substitute for real-proof
timings on matched inputs. The archived 1024-row proof control provides a
completed larger case, but its differing workload prevents a common scaling
curve with the new pilots. Results from one laptop also do not determine
server throughput or parallel scaling.

\subsection{Construct Validity}
The primary statement concerns the sign of exact quantized nearest-center
cost change. Ground-truth classification accuracy, quality before rounding,
encoder correctness and optimality of the proposed centers are different
constructs. Coverage before Full must be distinguished from the final
correctness obtained after paid exact work.

The ZKML control uses a published-primitive adapter with a different
commitment/backend and a stronger exact-cost output; it is not the complete
unmodified compiler. The PCA control shares our certificate wrapper.
These distinctions restrict attribution of runtime or proof-size differences.
The new paper-derived controls are independent reference implementations
at fixed empirical budgets, not reproductions of every optimization or
theorem parameter of their source systems. RaBitQ uses exact old-group
anchors and unauthenticated old cost, and its float64 statistical intervals
lack an outward-rounding proof. Native numerical time, statistical interval
resolution and deterministic cryptographic service time remain distinct.
Public routing and correction indices add observable outputs, while
authorized aggregate probes can reveal group statistics. No adaptive
assignment-privacy or record non-reconstruction guarantee is inferred.

